\section{Medidas de Posição}\label{sec:medidas de posicao}

Um grau maior de resumo dos dados pode ser obtido ao se computar determinadas estatísticas
numéricas que sejam representativas de todo o conjunto, e não somente de agrupamentos
deles, como é o caso das frequências absolutas e relativas.
Para tanto, podemos utilizar algumas das medidas que explicam uma determinada distribuição
vistas na parte I, geralmente usadas em conjunto para evitar uma redução muito drásticas
dos dados.
Nesta seção, veremos como obter, a partir da amostra, estimativas para medidas de posição,
em especial a moda, mediana e o valor esperado.

A primeira da qual iremos tratar será a moda, que pode ser interpretada como o valor mais
frequente numa distribuição.
\begin{def:moda}
Se $X$ é uma variável aleatória, e $f(x)$ é sua função de probabilidade, seja de massa ou
densidade, então a moda de $X$, denotada por $\mo{X}$, é o valor que satisfaz
\[
    \mo{X} = \argmax_x f(x).
\]
\end{def:moda}
Dada uma amostra $\{x_i\}_{i=1}^n$ de realizações para um atributo $X$, a
\textbf{moda amostral} -- estimativa da moda teórica a partir de uma amostra -- pode ser
definida simplesmente como o valor mais frequente nessa amostra.
A moda é adequada para atributos qualitativos e quantitativos discretos, ao passo que para
quantitativos contínuos ela pode ser pouco representativa devido ao grande número de
possibilidades de realizações.

A mediana é outra medida de centralidade muito utilizada em resumos, sendo definida a
seguir.
\begin{def:mediana}
Se $X$ é uma variável aleatória com função de distribuição $F(x)$, então a mediana,
denotada por $\med{X}$, satisfaz
\[
    F(\med{X}) = 0.5
\]
\end{def:mediana}
Com isso, a mediana é o valor que acumula metade da distribuição de probabilidade
de uma distribuição e, por consequência, sua contraparte amostral pode ser estimada ao
tomar a relização que tem frequência acumulada igual a $50\%$.
Se $N(x)$ é o número de observações tais que $X \le x$, então a mediana amostral é
o valor $x$ que satisfaz
\[
    \frac{N(x)}{n} = 0.5
\]
Computacionalemente, uma forma de obtê-la é ordenado as $n$ observações por seus valores,
e, caso $n$ seja ímpar, a mediana é o valor no índice $\frac{n}{2}$, ao passo que,
caso $n$ seja par, ela será a média aritmética entre as realizações de índice $\frac{n}{2}$
e $\frac{n+1}{2}$.
Há outros métodos computacionais para obter a mediana sem necessariamente ordenar os
dados, a exemplo do algoritmo \texttt{quickselect}.

Por fim, a média é a estimativa do valor esperado de uma variável $X$, sendo, pois,
o valor que corresponde ao centro de massa da amostra.
Para a obter, basta computar
\[
	\bar x = \frac{1}{n} \sum_{i=1}^n x_i,
\]
Tanto a média quanto a mediana são apropriadas apenas à atributos quantitativos, com
uma possível execeção àqueles que denotam uma ordenação, cuja média pode não ter uma
interpretação no contexto dos dados.
Caso estejamos lidando com intervalos discretizados desses valores, então a forma de
computar essas medidas é por estimativa, como, por exemplo, a tomar o ponto médio de
cada intervalo e usá-los para obter a média final.
